\lecture{13}{16. Oktober 2025}{Moment of Inertia}

\begin{problemwithimage}{f13_1}{F13-1}
  Determine the second moment of area about the $x$ axis of the cross-section of the I-beam.
  \textit{Remark}: Regions located symmetrically about the $x$-axis provide the same contribution to $I_x$.
\end{problemwithimage}

\begin{derivation}
  We know that the second moment of area of a rectangle is
  \begin{equation*}
    \overline{I}_{x'} = \frac{1}{12} bh^3
  .\end{equation*}
  For each of the three parts this corresponds to:
  \begin{align*}
    \overline{I}_{\mathrm{flange}_{x'}} &= \frac{1}{12} \cdot \qty{200}{\mm} \cdot \left( \qty{50}{\mm}  \right)^3  = \qty{2.0833e6}{\mm^4} \\
    I_{\mathrm{mid}_{x'}} &= \frac{1}{12} \cdot \qty{50}{\mm} \cdot \left( \qty{300}{\mm}  \right)^3  = \qty{1.125e8}{\mm^4}
  .\end{align*}
  To use the parallel axis theorem on the top and bottom flanges we must know their cross sectional areas:
  \begin{equation*}
    A_{\mathrm{flange}} = \qty{50}{\mm} \cdot \qty{200}{\mm}  = \qty{100}{\cm^2}
  .\end{equation*}
  Hence, using the parallel axis theorem
  \begin{align*}
    I_{\mathrm{flange}_{x'}} &= \overline{I}_{\mathrm{flange}_{x'}} + d_y^2 A_{\mathrm{flange}} \\
    &= \qty{2.0833e6}{\mm^{4}} + \left( \qty{175}{\mm}  \right)^2 \cdot \qty{100}{\cm^2}  \\
    &= \qty{308,333e6}{\mm^{4}} 
  .\end{align*}
  
  Hence
  \begin{align*}
    I_{\mathrm{tot}_{x'}} &= I_{\mathrm{mid}_{x'}} + 2 \cdot I_{\mathrm{flange}_{x'}} \\
    &= \qty{1.125e8}{\mm^{4}} + 2 \cdot \qty{308,333e6}{\mm^{4}}  \\
    &= \qty{729e6}{\mm^{4}} 
  .\end{align*}
\end{derivation}

\finalresult{I_{\mathrm{tot}_{x'}} = \qty{729e6}{\mm^{4}}}

\begin{problemwithimage}{f6_59}{6-59}
  Determine the moment of inertia for the shaded area about the $y$ axis.
\end{problemwithimage}

\begin{derivation}
  The formula for the second moment of area is
  \begin{equation*}
    I_y = \int_A x^2 \, \mathrm{d}A
  .\end{equation*}
  We therefore need to find an expression for our differential element $\mathrm{d}A$.

  We choose to use vertical strips with thickness $\mathrm{d}x$.
  At a given $x$, the shaded height is
  \begin{equation*}
    h(x) = 8 - \frac{1}{8}x^3
  .\end{equation*}
  So
  \begin{equation*}
    \mathrm{d}A = h(x) \, \mathrm{d}x = \left( 8 - \frac{1}{8} x^3 \right) \, \mathrm{d}x
  .\end{equation*}
  
  Our integral is then
  \begin{equation*}
    I_y = \int_{0}^{4} x^2 \left( 8 - \frac{1}{8} x^3 \right) \, \mathrm{d}x  = \num{85,33} 
  .\end{equation*}
  As our units are all in \unit{\m} this must be with dimension \unit{\m^4}.
\end{derivation}

\finalresult{I_y = \qty{85,33}{\m^4}}
